\documentclass[12 pt, letterpaper]{hmcpset}
\usepackage{hw-preamble}

\name{Jayden Thadani}
\class{MATH 139 --- Fourier Analysis}
\assignment{Homework \#3}
\duedate{14th February 2025}

\begin{document}

\begin{problem}[S \& S \S 2.6 \#6.]
	Let $f$ be the function defined on $[-\pi, \pi]$ by $f(\theta) = \lvert \theta \rvert$.
	\begin{enumerate}
		\item Draw the graph of $f$.
		\item Calculate the Fourier coefficients of $f$, and show that
			\[
				\hat{f}(n) = \begin{cases}
					\pi / 2 & \text{if } n = 0 \\
					\frac{-1 + (-1)^{n}}{\pi n^{2}} & \text{otherwise}.
				\end{cases}
			\]
		\item What is the Fourier series of $f$ in terms of sines and cosines?
		\item Taking $\theta = 0$, prove that
			\[
				\sum_{n \text{ odd} \ge 1} \frac{1}{n^{2}} = \frac{\pi^{2}}{8} \quad \text{and} \quad \sum_{n = 1}^{\infty} \frac{1}{n^{2}} = \frac{\pi^{2}}{6}
			\]
	\end{enumerate}
\end{problem}

\begin{solution}
	\begin{enumerate}
		\item The plot below was created in TikZ.
			\begin{figure}[H]
				\centering
				\begin{tikzpicture}
					\draw[->] (-4, 0) -- (4, 0) node[right] {$\theta$};
					\draw[->] (0, -1) -- (0, 4) node[above] {$f(\theta)$};
					\draw[scale = 1, domain = 0:3.14, smooth, variable = \x, blue] plot ({\x}, {\x});
					\draw[scale = 1, domain = -3.14:0, smooth, variable = \x, blue] plot ({\x}, {-\x});
				\end{tikzpicture}
				\caption{The absolute value function $f$ in blue}
			\end{figure}

		\item Notice that
			\[
				\hat{f}(0) = \frac{1}{2 \pi} \int_{-\pi}^{\pi} f(\theta) \, d \theta = \frac{1}{2 \pi} 2 \frac{\pi^{2}}{2} = \frac{\pi}{2}
			\]
			since the integral is just the area of two right angles each with base and height $\pi$.

			For other $n$, we can integrate by parts
			\begin{align*}
				\hat{f}(n) & = \frac{1}{2 \pi} \int_{-\pi}^{\pi} f(\theta) e^{-i n \theta} \, d\theta \\
					   & = \frac{1}{2 \pi} \left ( \int_{0}^{\pi} \theta e^{-i n \theta} \, d\theta - \int_{-\pi}^{0} \theta e^{-i n \theta} \, d\theta \right ) \\
					   & = \frac{1}{2 \pi} \left ( \left ( \frac{e^{-i n \theta}}{n^{2}} - \frac{\theta e^{-i n \theta}}{i n} \right ) \Big |_{0}^{\pi} - \left ( \frac{e^{-i n \theta}}{n^{2}} - \frac{\theta e^{-i n \theta}}{i n} \right ) \Big |_{-\pi}^{0} \right ) \\
					   & = \frac{1}{2 \pi} \left ( \left ( \frac{(-1)^{n}}{n^{2}} - \frac{\pi (-1)^{n}}{in} - \frac{1}{n^{2}} \right ) - \left ( \frac{1}{n^{2}} - \frac{\pi (-1)^{n}}{in} - \frac{(-1)^{n}}{n^{2}} \right ) \right ) \\
					   & = \frac{2}{2 \pi} \left ( \frac{(-1)^{n} - 1}{n^{2}}  \right ) \\
					   & = \frac{-1 + (-1)^{n}}{\pi n^{2}}.
			\end{align*}

		\item
			\[
				\boxed{
					f(\theta) \sim \frac{\pi}{2} - \sum_{n \text{ odd} \ge 1} \frac{4}{\pi n^{2}} \cos{(n \theta)}.
				}
			\]

			Notice that $\hat{f}(n) = \hat{f}(-n)$. Since $e^{i n \theta} + e^{-i n \theta} = 2 \cos{(n \theta)}$. Further since $\hat{f}(n) = 0$ for even $n$, we have
			\begin{align*}
				\sum_{n = -\infty}^{\infty} \hat{f}(n) e^{-i n \theta} & = \frac{\pi}{2} + \sum_{n = 1}^{\infty} 2 \hat{f}(n) \cos(n \theta) \\
										       & = \frac{\pi}{2} - \sum_{n \text{ odd} \ge 1} \frac{4}{\pi n^{2}} \cos{(n \theta)}.
			\end{align*}

		\item Since $f$ is continuous (at $0$, but also everywhere) and $\hat{f}(n) = O(1 / n^{2})$, the Fourier series recovers the function at $0$. Specifically
			\[
				\sum_{n = -\infty}^{\infty} \hat{f}(n) e^{i n 0} = f(0).
			\]
			
			Since $f(0) = 0$, evaluating the Fourier series at $0$ gives
			\[
				0 = \frac{\pi}{2} - \sum_{n \text{ odd} \ge 1} \frac{4}{\pi n^{2}} \cos{0}.
			\]
			Rearranging gives us
			\[
				\sum_{n \text{ odd} \ge 1} \frac{1}{n^{2}} = \frac{\pi^{2}}{8}.
			\]

			Every even number $n$ is $2^{k}$ times a unique odd number. Thus,
			\begin{align*}
				\sum_{n \text{ even} \ge 1} \frac{1}{n^{2}} & = \sum_{k = 1}^{\infty} \sum_{n \text{ odd} \ge 1} \frac{1}{(2^{k} n)^{2}} \\
									    & =  \sum_{k = 1}^{\infty} \frac{1}{4^k} \frac{\pi^{2}}{8} \\
									    & = \frac{\pi^{2}}{8} \frac{1 /4}{1 - 1/4} \\
									    & = \frac{\pi^{2}}{24}.
			\end{align*}

			It follows that
			\[
				\sum_{n = 1}^{\infty} \frac{1}{n^{2}} = \frac{\pi^{2}}{8} + \frac{\pi^{2}}{24} = \frac{\pi^{2}}{6}.
			\]
	\end{enumerate}
\end{solution}

\pagebreak

\begin{problem}[S \& S \S 2.6 \#9.]
	Let $f = \chi_{[a, b]}$ be the characteristic function of the interval $[a, b] \subseteq [-\pi, \pi]$, that is,
	\[
		\chi_{[a, b]}(x) = \begin{cases}
			1 & \text{if } x \in [a, b] \\
			0 & \text{otherwise}.
		\end{cases}
	\]
	\begin{enumerate}
		\item Show that the Fourier series of $f$ is given by
			\[
				f(x) \sim \frac{b - a}{2 \pi} + \sum_{n \neq 0} \frac{e^{-i n a}  - e^{-i n b}}{2 \pi i n} e^{i n x}
			\]
			The sum extends over all positive and negative integers excluding $0$.
		\item Show that if $a \neq -\pi$ or $b \neq -\pi$ and $a \neq b$, then the Fourier series does not converge absolutely for any $x$.
		\item However, prove that the Fourier series converges at every point $x$. What happens if $a = - \pi$ and $b = \pi$?
	\end{enumerate}
\end{problem}

\begin{solution}
	\begin{enumerate}
		\item Since $f = \chi_{[a, b]}$ has support only in $[a, b]$ where it is the identity, we only have to integrate the characters over that interval. Specifically, for all $n \neq 0$ 
			\begin{align*}
				\hat{f}(n) & = \frac{1}{2 \pi} \int_{a}^{b} e^{-i n x} \, dx \\
					   & = \frac{1}{2 \pi} \left ( \frac{e^{-i n x}}{-in} \right ) \Big |_{a}^{b} \\
					   & = \frac{e^{-i n a} - e^{-i n b}}{2 \pi i n}.
			\end{align*}
			For $n = 0$, integrating the character $e^{i 0 x} = 1$ over $[a, b]$ just gives
			\[
				\hat{f}(0) = \frac{1}{2 \pi} \int_{a}^{b}  \, dx = \frac{b - a}{2 \pi}.
			\]

		\item The magnitude of a generic term of the Fourier series is
			\[
				\frac{\lvert e^{-i n a} - e^{-i n b} \rvert}{2 \pi \lvert n \rvert} = \frac{\lvert e^{-i n a/2} e^{i n b/2} \rvert \lvert e^{i n (b - a)/2} - e^{-i n (b - a) / 2} \rvert}{2 \pi \lvert n \rvert} = \frac{\lvert \sin{(n \theta)} \rvert}{\pi \lvert n \rvert}.
			\]
			for $\theta = (b - a) / 2$. We claim there are infinitely many $n \theta$ such that $\lvert \sin{(n \theta)} \rvert$ is bounded away from $0$.

			We can see this by considering $\lvert \sin{(n \theta)} \rvert + \lvert \sin{((n + 1) \theta)} \rvert$. If $\theta$ is not a multiple of $\pi$, the terms are never simultaneously zero. Considering the continuous function on the compact set $[-\pi, \pi]$ (since $\sin$ is $2 \pi$ periodic), the function achieves a minimum away from $0$, and thus, can be bounded above $0$. That is, there are infinitely many $n$ such that $\lvert \sin{(n \theta)} \rvert \ge c > 0$. By comparison with the harmonic series the series diverges.

		\item We can use the Dirichlet test to show that the series converges for each $x$. Clearly $a_{n} = 1 / 2 \pi i n$ is monotonically decreasing. Further $b_{n} = (e^{-i n a} - e^{-i n b}) e^{i n x}$. Has bounded partial sums
			\[
				B_{N} = \sum_{1 \le \lvert n \rvert \le N} e^{-i n a} e^{i n x} - \sum_{1 \le \lvert n \rvert \le N} e^{-i n b} e^{i n x}.
			\]
			These partial sums are bounded since the Dirichlet kernels $D_{N} = \sum_{\lvert n \rvert \le N} e_{i n x}$ are bounded.

			If $a = -\pi$ and $b = \pi$, then for all $n \neq 0$,
			\[
				\hat{f}(n) = \frac{e^{-i n (-\pi)} - e^{-i n \pi}}{2 \pi i n} = \frac{-1 - (-1)}{2 \pi i n} = 0.
			\]
			but
			\[
				\hat{f}(0) = \frac{\pi - (-\pi)}{2 \pi} = 1.
			\]
			Thus, $\chi_{[a, b]}(x) = \hat{f}(0) = 1$ for all $x \in [-\pi, \pi]$.
	\end{enumerate}
\end{solution}

\pagebreak

\begin{problem}[S \& S \S 2.6 \#11.]
	Suppose that $\{ f_{k} \}_{k \in \mathbb{N}}$ is a sequence of Riemann integrable functions on the interval $[0, 1]$ such that
	\[
		\int_{0}^{1} \lvert f_{k}(x) - f(x) \rvert \, dx 
	\]
	as $k \to \infty$.

	Show that $\hat{f}_{k}(n) \to \hat{f}(n)$ uniformly in $n$ as $k \to \infty$.
\end{problem}

\begin{solution}
	We can compute the difference between $\hat{f}_{k}$ and $\hat{f}$ and bound it by pulling out absolute values. Specifically, we write
	\begin{align*}
		\lvert \hat{f}(n) - \hat{f}_{k}(n) \rvert & = \left \lvert \int_{-\pi}^{\pi} (f(x) - f_{k}(x)) e^{-i n x} \, dx  \right \rvert \\
							  & \le \lvert e^{-i n x} \rvert \left \lvert \int_{-\pi}^{\pi} \lvert f(x) - f_{k}(x) \rvert \, dx  \right \rvert \\
							  & = \left \lvert \int_{-\pi}^{\pi} \lvert f(x) - f_{k}(x) \rvert \, dx  \right \rvert.
	\end{align*}

	Thus, $\lvert \hat{f}(n) - \hat{f}_{k}(n) \rvert \to 0$ as well. Since our bound was independent of  $n$ ($\lvert e^{-i n x} \rvert = 1$ for all $n$) the convergence is uniform.
\end{solution}

\pagebreak

\begin{problem}[S \& S \S 2.6 \#12.]
	Prove that if a series of complex numbers $\sum c_{n}$ converges to $s$, then $\sum c_{n}$ is also Ces\`aro summable to $s$.
\end{problem}

\begin{solution}
	Write $s_{n}$ for the $n$th partial sum $\sum_{k = 1}^{n} c_{k}$ and $\sigma_{n}$ for the $n$th Ces\`aro sum $(\sum_{k = 1}^{n} s_{n})/n$

	Since $s_{n} \to s$, for each $\varepsilon > 0$, there is some $N \in \mathbb{N}$ such that all $n > N$ give $\lvert s - s_{n} \rvert < \varepsilon$. Suppose $\sum_{n = 1}^{N} \lvert s - s_{n} \rvert = S$. Choose $M \in \mathbb{N}$ large enough that $S / M < \varepsilon/2$ and $M > N$. Since all $\lvert s - s_{n} \rvert$ with $n > N$ have $s_{n} < \varepsilon / 2$, for any $m > M$ we have
	\[
		\frac{\lvert s - s_{N + 1} \rvert + \dots + \lvert s - s_{m} \rvert}{m} \le \frac{ \lvert s - s_{N + 1} \rvert + \dots + \lvert s - s_{m} \rvert}{m - N} < \varepsilon / 2.
	\]

	Thus, for all $m > M$
	\begin{align*}
		\lvert s - \sigma_{m} \rvert & = \left \lvert \frac{ms - \sum_{n = 1}^{m} s_{n}}{m} \right \rvert \\
					     & \le \frac{\sum_{n = 1}^{m} \lvert s - s_{n} \rvert}{m} \\
					     & = \frac{S}{m} + \frac{\sum_{n = N + 1}^{m} \lvert s - s_{n} \rvert}{m} \\
					     & < \frac{S}{M} + \frac{\varepsilon}{2} \\
					     & = \varepsilon.
	\end{align*}

	That is, we have a crude estimate showing $\sigma_{n} \to s$ as desired.
\end{solution}

\pagebreak

\begin{problem}[S \& S \S 2.6 \#15]
	Prove that the Fej\'er kernel is given by
	\[
		F_{N}(x) = \frac{1}{N} \frac{\sin^{2}{(Nx / 2)}}{\sin^{2}{(x / 2)}}.
	\]
\end{problem}

\begin{solution}
	Recall that the $k$th Dirichlet kernel is given by $D_{k}(x) = \sum_{\lvert n \rvert \le k} e^{i n x}$ and the $N$th Fej\'er kernel is $\frac{1}{N} \sum_{k = 0}^{N - 1} D_{k}$.

	Understanding each $D_{k}$ as a sum of geometric sums of $\omega$ and $\overline{\omega} = 1 / \omega$ for $\omega = e^{i x}$, we get
	\[
		D_{k} = \frac{\omega^{-k} - \omega^{k + 1}}{1 - \omega}.
	\]
	Thus,
	\[
		F_{N}(x) = \frac{1}{N} \frac{\sum_{k = 0}^{N - 1} \omega^{-k} - \omega^{k + 1}}{1 - \omega}.
	\]

	Notice that
	\begin{align*}
		4 \sin^{2}{(x / 2)} & = - (2 i \sin{(x / 2)})^{2} \\
				    & = - (e^{i x / 2} - e^{-i x / 2})^{2} \\
				    & = - (e^{i x} + e^{-i x} - 2 e^{i x / 2} e^{-i x / 2}) \\
				    & = (1 - \omega) + (1 - \omega^{-1}) \\
				    & = (1 - \omega) - \omega^{-1}(1 - \omega) \\
				    & = (1 - \omega)(1 - \omega^{-1})
	\end{align*}
	This motivates considering $(1 - \omega^{-1})$ times the numerator of the F\'ejer kernel
	\begin{align*}
		(1 - \omega^{-1}) \sum_{k = 0}^{N - 1} \omega^{-k} - \omega^{k + 1} & = \sum_{k = 0}^{N - 1} \omega^{-k} + \omega^{k} - \omega^{k + 1} - \omega^{-(k + 1)} \\
										    & = \omega^{0} + \omega^{0} - (\omega^{N} + \omega^{-N}) \\
										    & = 2 - 2 \cos{(N x)} \\
										    & = 2(1 - \cos{(N x)}) \\
										    & = 4 \sin^{2}{(N x) / 2}.
	\end{align*}
	Here the second equality is because the preceding sum is telescoping.

	Taking the quotient gives us the desired formula ---
	\[
		F_{N}(x) = \frac{1}{N} \frac{(1 - \omega^{-1}) \sum_{k = 0}^{N - 1} \omega^{-k} - \omega^{k + 1}}{(1 - \omega^{-1}) (1 - \omega)} = \frac{1}{N} \frac{\sin^{2}{(N x / 2)}}{\sin^{2}{(x / 2)}}.
	\]
\end{solution}

\pagebreak

\begin{problem}[S \& S \S 2.6 \#16]
	\begin{theorem}[Weierstrass approximation theorem]
		Let $f$ be a continuous function on the closed and bounded interval $[a, b] \subseteq \mathbb{R}$. Then, for any $\varepsilon > 0$, there exists a polynomial $P$ such that
		\[
			\sup_{x \in [a, b]} \lvert f(x) - P(x) \rvert < \varepsilon.
		\]
	\end{theorem}
	Prove this by applying Fej\'er's theorem and using the fact that the exponential function $e^{i x}$ can be approximated by polynomials uniformly on any interval.
\end{problem}

\begin{solution}
	Recall that uniform convergence implies $L^{\infty}$ convergence.

	By F\'ejer's theorem, there exists a trigonometric polynomial $p$ such that
	\[
		\sup_{x \in [a, b]} \lvert f(x) - p(x) \rvert < \varepsilon /2.
	\]
	We are given that there exists a polynomial $P$ such that
	\[
		\sup_{x \in [a, b]} \lvert p(x) - P(x) \rvert < \varepsilon / 2.
	\]
	Thus, by the triangle inequality,
	\[
		\sup_{x \in [a, b]} \lvert f(x) - P(x) \rvert < \varepsilon.
	\]
\end{solution}

\end{document}
